% ============================================================
%  Fire Detection — Algorithm Evaluation Report
%  Smart Thermal System for Patient Safety Monitoring
%  Afeka College of Engineering, Tel-Aviv
%
%  Compiled with: pdflatex (TeX Live / MiKTeX / Overleaf)
% ============================================================
\documentclass[11pt, a4paper]{article}

% ---- Packages -----------------------------------------------
\usepackage[a4paper, margin=2.5cm]{geometry}
\usepackage[T1]{fontenc}
\usepackage{lmodern}
\usepackage{booktabs}
\usepackage{multirow}
\usepackage{array}
\usepackage{amsmath}
\usepackage{amssymb}
\usepackage[table]{xcolor}   % loads colortbl for \cellcolor
\usepackage{graphicx}
\usepackage{caption}
\usepackage{hyperref}
\usepackage{parskip}
\usepackage{enumitem}
\usepackage{float}

% Fallback: some MiKTeX builds don't expose these array.sty internals
% (\do@row@strut is invoked by every p/m/b column). Define safe stubs only if
% missing so p-columns typeset without raising "Undefined control sequence".
\providecommand{\arraybackslash}{\let\\\tabularnewline}
\makeatletter
\@ifundefined{do@row@strut}{\let\do@row@strut\relax}{}
\makeatother

% ---- Colours ------------------------------------------------
\definecolor{tpgreen}{RGB}{0,128,0}
\definecolor{tnblue}{RGB}{31,73,125}
\definecolor{fpred}{RGB}{192,0,0}
\definecolor{fnyellow}{RGB}{150,100,0}
\definecolor{headerblue}{RGB}{31,73,125}
\definecolor{cellgreenlight}{RGB}{198,239,206}
\definecolor{cellbluelight}{RGB}{221,235,247}
\definecolor{cellredlight}{RGB}{255,199,206}
\definecolor{cellyellowlight}{RGB}{255,235,156}

% ---- Hyperref setup -----------------------------------------
\hypersetup{
    colorlinks = true,
    linkcolor  = headerblue,
    urlcolor   = headerblue,
    citecolor  = headerblue
}

% ---- Confusion matrix macro ---------------------------------
% Usage: \confmat{TP}{FN}{FP}{TN}{pos_total}{neg_total}
\newcommand{\confmat}[6]{%
    \def\cmTP{#1}\def\cmFN{#2}\def\cmFP{#3}\def\cmTN{#4}%
    \def\cmPos{#5}\def\cmNeg{#6}%
    \vspace{4pt}
    \begin{center}
    \renewcommand{\arraystretch}{1.8}%
    \begin{tabular}{lcc}
        \hline
        \rule{0pt}{2.4ex}%
        & \textbf{Pred:\ Fire} & \textbf{Pred:\ Safe} \\[2pt]
        \hline
        \rule{0pt}{2.6ex}%
        \textbf{Truth:\ Fire} (\cmPos\ frames)
            & \cellcolor{cellgreenlight}\textcolor{tpgreen}{\textbf{TP}}\;$=$\;\cmTP
            & \cellcolor{cellyellowlight}\textcolor{fnyellow}{\textbf{FN}}\;$=$\;\cmFN \\[4pt]
        \textbf{Truth:\ Safe} (\cmNeg\ frames)
            & \cellcolor{cellredlight}\textcolor{fpred}{\textbf{FP}}\;$=$\;\cmFP
            & \cellcolor{cellbluelight}\textcolor{tnblue}{\textbf{TN}}\;$=$\;\cmTN \\[2pt]
        \hline
    \end{tabular}
    \end{center}
    \vspace{4pt}
}

% ---- Title --------------------------------------------------
\title{%
    {\large Afeka College of Engineering, Tel-Aviv}\\[0.4em]
    {\large Smart Thermal System for Patient Safety Monitoring}\\[0.8em]
    {\LARGE\bfseries Fire Detection — Algorithm Evaluation Report}\\[0.4em]
    {\normalsize MLX90640 Thermal Sensor \quad 22 Scenes \quad
     5{,}943 Annotated Frames}
}
\author{%
    Guy Chen \and Yaniv Blau \and Roy Lieberman\\[0.3em]
    \small Supervisor: Or Zilberberg \qquad Advisor: Dr.\@ Oshrit Hoffer
}
\date{June 2026}

% =============================================================
\begin{document}
\maketitle
\tableofcontents
\newpage

% =============================================================
\section{Overview}
% =============================================================

This report evaluates the fire detection algorithms of the Smart Thermal System
(\S4.4.4 of the engineering report) on the full annotated dataset.  Three detector
configurations are compared:

\begin{enumerate}[leftmargin=2em]
    \item \textbf{Otsu Pipelined (original)} — the rule-based three-stage detector
          using Otsu segmentation plus morphological erosion/dilation, as originally
          specified.
    \item \textbf{Otsu Hot-Pixel Bypass (improved)} — the same detector after a
          targeted fix that segments directly on an absolute temperature threshold and
          removes the erosion step (Section~\ref{sec:bypass}).
    \item \textbf{Fire SVM} — the feature-based RBF support-vector classifier.
\end{enumerate}

Every frame is reduced to a binary verdict — \textbf{Fire} (an alert of any level)
or \textbf{Safe} — and scored with a confusion matrix and the four derived metrics
(accuracy, precision, recall, F1).  The central finding is that fire detection on this
sensor is fundamentally limited by the \emph{thermal signature} of the labelled fires,
not only by algorithm design: the headline numbers must be read together with the
data analysis in Section~\ref{sec:ceiling}.

% =============================================================
\section{Dataset and Experimental Setup}
% =============================================================

\subsection{Dataset}

Recordings were collected at the Beer Sheva Mental Health Center using three MLX90640
sensors (32$\times$24 px, 8\,Hz).  The evaluation uses \textbf{22 scenes} totalling
\textbf{5{,}943 annotated frames} across all three channels.  Fire is annotated as
YOLO class~0; a frame is a fire \emph{positive} if it carries at least one fire box.

\begin{table}[H]
\centering
\caption{Fire-positive vs.\ fire-negative frame counts (all channels).}
\renewcommand{\arraystretch}{1.3}
\begin{tabular}{lcc}
\toprule
\textbf{Scene group} & \textbf{Fire frames (Pos)} & \textbf{Safe frames (Neg)} \\
\midrule
\texttt{3pplfire}                 & 42 & 213 \\
\texttt{onepersonfire}            & 31 & 230 \\
\texttt{setup2\_fire\_lighter}    &  5 & 301 \\
\addlinespace
Hot distractors (hairdryer, PC screen) & 0 & 528 \\
Other 17 scenes (people, motion)       & 0 & 4{,}393 \\
\midrule
\textbf{Total} & \textbf{78} & \textbf{5{,}865} \\
\bottomrule
\end{tabular}
\label{tab:counts}
\end{table}

The dataset is \textbf{extremely imbalanced}: only 78 of 5{,}943 frames
($\approx 1.3\,\%$) contain fire, spread across just three scenes.  Two scenes are
deliberately retained as \emph{hard negatives}: \texttt{2pplwithhairdryer} (a hot air
jet) and \texttt{emptyroomwithpcscreen} (a warm monitor).  These contain genuinely hot
regions but no fire, and are the most demanding test of precision.

\subsection{Evaluation Protocol}

The two \textbf{Otsu} configurations are rule-based and require no training, so they
are evaluated on the \textbf{full dataset} (all 5{,}943 frames).  The \textbf{Fire SVM}
is trainable and is evaluated on a held-out \textbf{70\,/\,30 test split} stratified per
(scene, channel, label): 4{,}128 training frames (53 positive) and 1{,}815 test frames
(25 positive).  For a fair head-to-head, the improved Otsu detector is \emph{also}
reported on that same test split in Section~\ref{sec:headtohead}.  The detector's
temporal state is reset at every scene boundary.

% =============================================================
\section{Metric Definitions}
% =============================================================

\subsection{Confusion Matrix}

Each frame yields one of four outcomes:

\begin{itemize}[leftmargin=2em]
    \item \textbf{TP} (True Positive): fire present and correctly alarmed.
    \item \textbf{TN} (True Negative): no fire and system correctly silent.
    \item \textbf{FP} (False Positive): no fire but alarm raised — a \emph{false alarm}.
    \item \textbf{FN} (False Negative): fire present but missed — a \emph{missed fire}.
\end{itemize}

\subsection{Derived Metrics}

\begin{align}
    \text{Accuracy}  &= \frac{TP + TN}{TP + TN + FP + FN} \label{eq:acc}  \\[6pt]
    \text{Precision} &= \frac{TP}{TP + FP}                \label{eq:prec} \\[6pt]
    \text{Recall}    &= \frac{TP}{TP + FN}                \label{eq:rec}  \\[6pt]
    \text{F1}        &= \frac{2 \cdot \text{Precision} \cdot \text{Recall}}
                              {\text{Precision} + \text{Recall}}            \label{eq:f1}
\end{align}

\paragraph{Accuracy} (Eq.\,\ref{eq:acc}) is the fraction of all frames classified
correctly.  Under 1.3\,\% positives it is \textbf{badly misleading}: a detector that
never alarms scores $98.7\,\%$ accuracy while missing every fire.  Accuracy is reported
for completeness only and should not be used to rank detectors here.

\paragraph{Precision} (Eq.\,\ref{eq:prec}) is the fraction of fire alarms that are
genuine.  Low precision means frequent false alarms, which cause \emph{alert fatigue}
and lead staff to disable the system — especially damaging in a psychiatric ward.

\paragraph{Recall} (Eq.\,\ref{eq:rec}) is the fraction of real fires detected.  It is
the \textbf{primary safety metric}: a missed ignition is a direct, potentially
life-threatening system failure.

\paragraph{F1-Score} (Eq.\,\ref{eq:f1}) is the harmonic mean of precision and recall.
It is the headline metric because accuracy is uninformative under this imbalance.

% =============================================================
\section{Algorithms}
% =============================================================

\begin{description}[leftmargin=1.5em, style=nextline]
    \item[Otsu Pipelined (original) — rule-based, no training.]
        Variance gate $\to$ Otsu segmentation $\to$ 1$\times$ erosion + 2$\times$
        dilation $\to$ hierarchical $T_{\max}$/area rules $\to$ temporal mass-gradient
        growth check.
    \item[Otsu Hot-Pixel Bypass — rule-based, no training.]
        Same pipeline, but segmentation thresholds directly at $T_{\text{ign}}$ via
        connected components (pixel-count area, no erosion), preceded by dead-pixel
        despiking.  See Section~\ref{sec:bypass}.
    \item[Fire SVM — classical ML, RBF support-vector classifier.]
        Otsu region proposal $\to$ 6 scalar blob features ($T_{\max}$,
        $T_{\text{mean}}$, $T_{\text{std}}$, area, skewness, kurtosis) $\to$
        standardisation $\to$ RBF SVM decision.
\end{description}

Decision thresholds (both Otsu variants): $T_{\text{ign}}=45\,^\circ$C,
$T_{\text{fire}}=60\,^\circ$C, $A_{\text{limit}}=23$\,px ($\approx3\,\%$ of the frame).

% =============================================================
\section{Results — Otsu Detector (full dataset)}
% =============================================================

\subsection{Confusion Matrices}

\subsubsection*{Otsu Pipelined (original)}
\confmat{7}{71}{3}{5862}{78}{5865}

\subsubsection*{Otsu Hot-Pixel Bypass (improved)}
\confmat{21}{57}{31}{5834}{78}{5865}

\subsection{Per-Scene Breakdown}

\begin{table}[H]
\centering
\caption{Per-scene confusion counts on the three fire scenes and two hot distractors
(full dataset).  ``Orig.''\ = Otsu Pipelined; ``Bypass''\ = Hot-Pixel Bypass.}
\renewcommand{\arraystretch}{1.3}
\begin{tabular}{l cccc c cccc}
\toprule
& \multicolumn{4}{c}{\textbf{Original}} & & \multicolumn{4}{c}{\textbf{Bypass}} \\
\cmidrule(lr){2-5}\cmidrule(lr){7-10}
\textbf{Scene} & TP & TN & FP & FN & & TP & TN & FP & FN \\
\midrule
\texttt{3pplfire}              & 7 & 213 & 0 & 35 & & 20 & 213 & 0  & 22 \\
\texttt{onepersonfire}         & 0 & 230 & 0 & 31 & &  1 & 228 & 2  & 30 \\
\texttt{setup2\_fire\_lighter} & 0 & 301 & 0 &  5 & &  0 & 299 & 2  &  5 \\
\addlinespace
\texttt{2pplwithhairdryer}     & 0 & 298 & 2 &  0 & &  0 & 288 & 12 &  0 \\
\texttt{emptyroomwithpcscreen} & 0 & 228 & 0 &  0 & &  0 & 224 & 4  &  0 \\
\bottomrule
\end{tabular}
\label{tab:perscene}
\end{table}

\subsection{Metric Summary}

\begin{table}[H]
\centering
\caption{Otsu detector — full-dataset metrics (5{,}943 frames).}
\renewcommand{\arraystretch}{1.3}
\begin{tabular}{lcccccccc}
\toprule
\textbf{Configuration}
    & \textbf{Acc} & \textbf{Prec} & \textbf{Rec} & \textbf{F1}
    & \textbf{TP}  & \textbf{TN}   & \textbf{FP}  & \textbf{FN} \\
\midrule
Otsu Pipelined (original)
    & 98.8\% & 70.0\% & 9.0\% & \textbf{15.9\%}
    & 7  & 5862 &  3 & 71 \\
Otsu Hot-Pixel Bypass
    & 98.5\% & 40.4\% & \textbf{26.9\%} & \textbf{32.3\%}
    & 21 & 5834 & 31 & 57 \\
\bottomrule
\end{tabular}
\label{tab:otsu}
\end{table}

% =============================================================
\section{The Hot-Pixel Bypass Fix}\label{sec:bypass}
% =============================================================

The original detector's recall of $9.0\,\%$ was traced to two interacting defects in
the segmentation stage:

\begin{enumerate}[leftmargin=2em]
    \item \textbf{Erosion deletes single-pixel fires.}  At 32$\times$24 resolution a
          real fire is frequently a \emph{single} hot pixel.  In \texttt{onepersonfire}
          one frame contains a genuine $72.5\,^\circ$C pixel, yet the $1\times$ erosion
          (3$\times$3 kernel) removes any pixel lacking hot 8-connected neighbours,
          deleting the fire entirely before classification.
    \item \textbf{Otsu's threshold is anchored by warm bodies.}  With people in frame,
          Otsu selects a threshold around $25\,^\circ$C, so the dominant ``hot'' blob is
          a person ($\approx39\,^\circ$C) and the $72.5\,^\circ$C fire pixel is never
          isolated as its own region.
\end{enumerate}

The fix replaces Otsu+erosion in the alert path with direct thresholding at
$T_{\text{ign}}$ using connected-component labelling (\emph{pixel-count} area, so a lone
pixel has area~1, not~0), preceded by a despiking filter that neutralises isolated
dead-pixel glitches (single pixels reading $800$--$935\,^\circ$C).  The hierarchical
classifier already acted only on blobs with $T_{\max}\ge T_{\text{ign}}$, so this is
equivalent in intent while recovering the small fires erosion discarded.

\paragraph{Effect.}  Recall improved $\mathbf{9.0\,\% \rightarrow 26.9\,\%}$
($3\times$, TP $7\rightarrow21$) and F1 $\mathbf{15.9\,\% \rightarrow 32.3\,\%}$.
Localisation quality on \texttt{3pplfire} rose from a mean IoU of $1.0\,\%$ to
$17.3\,\%$.  The cost is precision: $70\,\% \rightarrow 40.4\,\%$ (FP $3\rightarrow31$),
because thresholding at $45\,^\circ$C now also flags warm distractor pixels — most
visibly the hairdryer jet (12 FP) and PC screen (4 FP).  This precision/recall trade-off
is examined in Sections~\ref{sec:ceiling} and~\ref{sec:interp}.

% =============================================================
\section{The Recall Ceiling — A Data Limitation}\label{sec:ceiling}
% =============================================================

The improved recall of $26.9\,\%$ is not an algorithmic shortfall; it is essentially the
\textbf{physical ceiling} imposed by the labels.  Analysing the 78 fire-positive frames:

\begin{table}[H]
\centering
\caption{Peak-temperature analysis of the 78 fire-labelled frames.}
\renewcommand{\arraystretch}{1.3}
\begin{tabular}{lc}
\toprule
\textbf{Property} & \textbf{Value} \\
\midrule
Median peak temperature of a fire frame      & $36.5\,^\circ$C \\
Minimum peak temperature                      & $28.7\,^\circ$C \\
Fire frames with any pixel $>45\,^\circ$C ($T_{\text{ign}}$) & 21 / 78 \;(26.9\,\%) \\
Fire frames with any pixel $>60\,^\circ$C ($T_{\text{fire}}$) & 9 / 78 \;(11.5\,\%) \\
\bottomrule
\end{tabular}
\label{tab:ceiling}
\end{table}

The median fire frame peaks at only $36.5\,^\circ$C — indistinguishable from a human body
($33$--$37\,^\circ$C) — and \textbf{73\,\% of fire-labelled frames contain no pixel above
the $45\,^\circ$C ignition threshold}.  The annotation boxes the fire \emph{object}
wherever it is visible (smouldering, distant, or sub-pixel), not only frames carrying a
hot thermal signature.  Consequently \textbf{no absolute-temperature detector can exceed
$\approx27\,\%$ recall} on this label set, and the bypass detector has now reached that
ceiling.  This is the single most important conclusion of the study: meaningful further
gains require re-annotation, not algorithm tuning (Section~\ref{sec:reco}).

% =============================================================
\section{Head-to-Head — Test Split}\label{sec:headtohead}
% =============================================================

To compare the trainable SVM against the rule-based detector on identical data, both are
evaluated on the same held-out test split (1{,}815 frames, 25 positive).

\subsection{Confusion Matrices}

\subsubsection*{Otsu Hot-Pixel Bypass}
\confmat{5}{20}{3}{1787}{25}{1790}

\subsubsection*{Fire SVM (RBF)}
\confmat{8}{17}{0}{1790}{25}{1790}

\subsection{Metric Summary}

\begin{table}[H]
\centering
\caption{Test-split metrics (1{,}815 frames, 25 positive).}
\renewcommand{\arraystretch}{1.3}
\begin{tabular}{lcccccccc}
\toprule
\textbf{Detector}
    & \textbf{Acc} & \textbf{Prec} & \textbf{Rec} & \textbf{F1}
    & \textbf{TP}  & \textbf{TN}   & \textbf{FP}  & \textbf{FN} \\
\midrule
Otsu Hot-Pixel Bypass
    & 98.7\% & 62.5\% & 20.0\% & \textbf{30.3\%}
    & 5 & 1787 & 3 & 20 \\
Fire SVM (RBF)
    & 99.1\% & \textbf{100.0\%} & \textbf{32.0\%} & \textbf{48.5\%}
    & 8 & 1790 & 0 & 17 \\
\bottomrule
\end{tabular}
\label{tab:headtohead}
\end{table}

% =============================================================
\section{Consolidated Comparison}
% =============================================================

\begin{table}[H]
\centering
\caption{Complete metric comparison across all three fire-detector configurations.}
\renewcommand{\arraystretch}{1.35}
\begin{tabular}{llcccc}
\toprule
\textbf{Detector} & \textbf{Eval set}
    & \textbf{Accuracy} & \textbf{Precision} & \textbf{Recall} & \textbf{F1} \\
\midrule
Otsu Pipelined (original)
    & Full (5{,}943) & 98.8\% & 70.0\%          & 9.0\%  & 15.9\% \\
Otsu Hot-Pixel Bypass
    & Full (5{,}943) & 98.5\% & 40.4\%          & 26.9\% & 32.3\% \\
\addlinespace
Otsu Hot-Pixel Bypass
    & Test (1{,}815) & 98.7\% & 62.5\%          & 20.0\% & 30.3\% \\
Fire SVM (RBF)
    & Test (1{,}815) & 99.1\% & \textbf{100.0\%}& \textbf{32.0\%} & \textbf{48.5\%} \\
\bottomrule
\end{tabular}
\label{tab:full}
\end{table}

% =============================================================
\section{Interpretation}\label{sec:interp}
% =============================================================

\paragraph{Fire SVM is the strongest detector.}  At F1\,=\,$48.5\,\%$ it leads every
configuration, and it is the only detector with \textbf{zero false alarms} —
crucially, it correctly stays silent on \emph{both} hot distractors (hairdryer and PC
screen).  It achieves this by classifying on blob \emph{shape} statistics (skewness,
kurtosis, $T_{\text{std}}$, area) rather than temperature alone: a compact, sharply
peaked fire signature is separable from a broad, smooth hot surface even when their peak
temperatures coincide.  On the common test split it beats the rule-based detector on
every metric ($+12$\,pp recall, $+37.5$\,pp precision, $+18.2$\,pp F1).

\paragraph{The Otsu bypass is a high-recall, low-precision gate.}  The fix tripled recall
and is the right behaviour for a rule-based first stage, but a pure temperature threshold
\emph{cannot} separate a $50\,^\circ$C fire pixel from a $50\,^\circ$C hairdryer pixel —
they are thermally identical.  Its false alarms arrive through the immediate
\textsc{Ignition-Source} path, which by design alarms instantly on any small hot blob and
therefore bypasses the temporal growth check that would otherwise reject static heat
sources.

\paragraph{Accuracy is meaningless here.}  All four rows in Table~\ref{tab:full} report
$98.5$--$99.1\,\%$ accuracy despite F1 ranging from $15.9\,\%$ to $48.5\,\%$.  This is the
imbalance artefact warned of in Eq.\,\ref{eq:acc}: with 98.7\,\% of frames negative, even
a near-blind detector looks ``accurate''.  F1 and recall are the only honest rankings.

\paragraph{Recall is capped by the labels, not the model.}  Both detectors plateau near
the $27\,\%$ ceiling of Section~\ref{sec:ceiling}.  The SVM slightly exceeds the naive
$>45\,^\circ$C threshold by learning warm-but-not-glowing fire textures, but the
$\approx73\,\%$ of sub-$45\,^\circ$C fire frames remain out of reach for any detector
operating on this sensor.

% =============================================================
\section{Conclusions and Recommendations}\label{sec:reco}
% =============================================================

\subsection{Findings}

\begin{enumerate}[leftmargin=2em]
    \item \textbf{Fire SVM} is the recommended primary detector: F1\,=\,$48.5\,\%$,
          precision\,=\,$100\,\%$, zero false alarms including on hot distractors.
    \item The \textbf{hot-pixel bypass} fixed a genuine segmentation bug, tripling the
          rule-based detector's recall ($9.0\to26.9\,\%$) and restoring usable
          localisation (IoU $1\to17\,\%$); it is best used as a high-recall pre-filter,
          not a standalone alarm.
    \item \textbf{Accuracy is not a valid metric} under this 1.3\,\% imbalance; F1 and
          recall govern all conclusions.
    \item Recall is bounded at $\approx27\,\%$ by a \textbf{label/physics mismatch}: most
          fire-labelled frames carry no above-threshold thermal signature.
\end{enumerate}

\subsection{Recommended Actions}

\begin{description}[leftmargin=2.5em, style=nextline]
    \item[1. Re-annotate fire by thermal signature (highest leverage).]
        Decide whether ``fire'' denotes a \emph{visible object} or a \emph{detectable
        thermal hazard}.  For a thermal-only safety system it must be the latter.  Until
        the sub-$45\,^\circ$C frames are relabelled or excluded, the $27\,\%$ recall
        ceiling is immovable regardless of algorithm.
    \item[2. Add temporal confirmation to the ignition path.]
        Extending the mass-gradient persistence check to small \textsc{Ignition-Source}
        blobs (not only large \textsc{Potential-Fire} blobs) should reject the static
        hairdryer/PC hot spots and recover much of the precision lost by the bypass,
        without sacrificing the recall gain.
    \item[3. Deploy the SVM with imbalance handling.]
        Retrain with class weighting or positive oversampling and add temporal/shape
        features to trade some of its perfect precision for higher recall, which is the
        priority safety metric.
    \item[4. Keep the despiking filter in preprocessing.]
        Dead-pixel glitches ($800$--$935\,^\circ$C) are a real false-alarm hazard for any
        $T_{\max}$ rule and must be neutralised before field deployment.
\end{description}

\end{document}
